% Written By Enoch Yu
% Downloaded from archive.enochyu.com

\documentclass[12pt]{article}
\usepackage{graphicx}
\usepackage{amsmath}
\usepackage{amssymb}
\usepackage{hyperref}
\hypersetup{
    colorlinks=true,
    linkcolor=blue,
    filecolor=magenta,      
    urlcolor=cyan,
    pdftitle={Problem by Enoch},
    pdfpagemode=FullScreen,
    }
\usepackage{fancyhdr}
\pagestyle{fancy}
\fancyhead[R]{Enoch Yu}

\title{Problem by Enoch}
\author{Enoch Yu}
\date{March 2025}

\begin{document}
\textbf{Problem} There are four magnificent cities named \textit{Happy}, \textit{Content}, \textit{Merry}, and \textit{Jolly}. One day, countable virus named \textit{MathVirus} has entered each city.
\\\\
The number of viruses that entered the \textit{Merry} city is the product of the number of virus that entered \textit{Happy} city and \textit{Content} city. Moreover, the sum of the number of virus that entered \textit{Happy} city and \textit{Content} city equals the number of virus that entered \textit{Jolly} city. In addition, the number of virus that entered \textit{Merry} city squared is equal to the product of 2025 and the number of virus that entered \textit{Jolly} city squared.
\\\\
A biologist in \textit{Happy} city is required to count the number of the \textit{MathVirus} that spread to \textit{Happy} city. Please help the biologist to find all possible number of virus that dispersed to \textit{Happy} city for further investigation! (The number of \textit{MathVirus} for each city does not change.)

\newpage

\noindent
\textbf{Solution}
\newline
Let\\
the number of virus in \textit{Happy} city $H$, \\
the number of virus in \textit{Content} city $C$, \\
the number of virus in \textit{Merry} city $M$, \\
the number of virus in \textit{Jolly} city $J$.
\\\\
$M$=$H$$C$
\newline
$J$=$H$+$C$
\newline
$M^2$=2025$J^2$
\newline
\newline
Because the number of the virus in each city is countable, the quantity of \textit{MathVirus} in each city is natural number. \\
$\therefore M=45J$
\newline
\newline
Because virus has entered each city, $J$ cannot be zero. \\
$\frac{M}{J}=45$
\newline
\newline
By substitution, \\
$\frac{HC}{H+C}=45$
\newline
$HC=45H+45C$
\newline
$HC-45H-45C=0$
\newline
$HC-45H-45C+2025=2025$
\newline
$(H-45)(C-45)=2025$
\newline
\newline
The possible values for $H-45$ are the factors of 2025.
\begin{align*}
    &H-45=1 \\
    &H-45=3 \\
    &H-45=5 \\
    &H-45=9 \\
    &H-45=15 \\
    &H-45=25 \\
    &H-45=27 \\
    &H-45=45 \\
    &H-45=75 \\
    &H-45=81 \\
    &H-45=135 \\
    &H-45=225 \\
    &H-45=405 \\
    &H-45=675 \\
    &H-45=2025
\end{align*}
\noindent
Therefore, the possible values of $H$ are 46, 48, 50, 54, 60, 70, 72, 90, 120, 126, 180, 270, 450, 720, 2070.
\\\\
\\\\
Conclusion: \textbf{The possible numbers of virus that entered to \textit{Happy} city are 46, 48, 50, 54, 60, 70, 72, 90, 120, 126, 180, 270, 450, 720, 2070}.
\end{document}

